\chapter{지식과 지혜 (글쓰기 첨삭)}

\begin{summarybox}{한눈에}
학생이 ``지식과 지혜'' 게시판에 올린 글을 \textbf{본 기관 학생 글만 골라 첨삭}합니다. 빨간 펜으로 직접 표시하거나, AI 자몽을 써서 자동 첨삭을 받을 수 있어요.
\end{summarybox}

\kfShot{11-wisdom}{지식과 지혜 — 본 기관 학생 글 목록}

\section*{화면 구조}

\begin{screenbox}{본 기관 학생 글 목록}
\noindent
\begin{tabular}{@{}p{2.5cm}p{2cm}p{2cm}p{3cm}p{2cm}p{2cm}@{}}
\toprule
\textbf{학생} & \textbf{레벨} & \textbf{주제} & \textbf{제출일} & \textbf{첨삭 여부} & \textbf{액션} \\
\midrule
홍길동 & 러셀1 & 환경 보호 주장 & 2026-05-04 & 미첨삭 & \inacttab{열기} \\
김영희 & 러셀2 & 미래 직업 상상 & 2026-05-05 & 완료 & \inacttab{열기} \\
\bottomrule
\end{tabular}
\end{screenbox}

\section*{여기서 할 수 있는 일}
\begin{itemize}[leftmargin=1.2em]
  \item 본 기관 학생이 쓴 글만 \textbf{필터된 목록} 으로 조회 (다른 기관 글 자동 제외)
  \item 글 클릭 → \textbf{원고지 그리드 + 빨간 펜 어노테이션}으로 직접 첨삭
  \item \textbf{AI 첨삭} 버튼 — 자몽 1개 차감 → AI 가 코멘트 + 교정 자동 생성 (몇 분)
  \item \textbf{AI 마크다운 OCR} — 학생이 종이로 쓴 사진을 업로드하면 텍스트로 변환 (자몽 1개)
  \item \textbf{일괄 AI 첨삭} — 여러 학생 글을 동시에 (각 자몽 차감)
  \item 첨삭 후 학생 화면에서 즉시 결과 확인
\end{itemize}

\section*{자주 쓰는 흐름}
\begin{enumerate}
  \item 학생이 글 제출 → 학습 계획표 ``글쓰기'' 탭 또는 이 화면에 신규 항목으로 표시
  \item 글 클릭 → 원고지 그리드 + 학생 본문 표시
  \item 빨간 펜으로 어노테이션 (단어 클릭, 의견 입력, 위치 표시)
  \item 또는 ``AI 첨삭'' 클릭 → 자몽 차감 안내 → 첨삭 결과 (몇 분 후)
  \item 저장 후 학생 화면에 결과 즉시 노출
\end{enumerate}

\begin{tipbox}{알아두면 좋은 점}
\begin{itemize}[leftmargin=1.2em]
  \item 본 기관 학생 글만 보입니다 — URL 직접 입력해도 다른 기관 글은 403
  \item AI 첨삭은 \textbf{Sonnet 4.6} 모델 — 자몽 1개 (200원 상당)
  \item OCR 은 학생이 종이에 손글씨로 쓴 사진을 텍스트로 변환할 때 — 자몽 1개
  \item 자몽 잔액 부족 시 자동으로 자몽 충전 페이지 안내가 뜹니다 (Chapter 12)
\end{itemize}
\end{tipbox}
